\part[When there are too few Euclidean directions]{When there are too few \titlebreak Euclidean directions}
\section{Field transfer and cube exclusion}\label{sec:prices}

We prove Theorem~\ref{thm:flats} by contradiction. If some Euclidean
dimension cannot be realized, we will find a finite field over which
spectral spaces uniformly exclude the corners of a fixed Euclidean cube.
This exclusion constrains how points can move relative to a maximal flat
in the routing argument of Section~\ref{sec:protocol}. Section~\ref{sec:ports}
uses the same exclusion to construct sparse graphs for compression.

\par\Needspace{4\baselineskip}
\begin{lemma}[Field transfer]\label{lem:transfer}
Fix an integer $a\ge1$. If spectral $a$-flats exist over every finite prime field, then
they exist over every field.
\end{lemma}
\begin{proof}
Fix any finite list in $\mathbb R^a$ and realize its distances in the assumed
spectral flat over each prime field. In characteristic $p$,
Lemma~\ref{lem:comparison} lets us extend scalars from $\mathbb F_p$
while preserving the evaluations on the entire list of pairwise binary wires. Using the original common scale gives the
prescribed distances. In characteristic zero, the same distance array occurs
over every prime field, so a limit argument
realizes it over $\mathbb Q$ and hence over the given field. Indeed, for a
positive-diameter array, \eqref{eq:span-price} bounds the scale below
by its diameter and above by the length of a spanning tree. Reduce the countably many rational tensors in
standard bases at primes avoiding their denominators, and pass to a common
subsequence on which scales and evaluations converge. Every rational
restriction and every tensor-product or direct-sum identity survives reduction
at all sufficiently large primes. The limits therefore define a spectral
point giving the required array. A singleton array is immediate.

Exhaust $\mathbb Q^a$ by finite lists containing the origin. Place each list in
a complete spectral space over the chosen field, based at the image of zero.
A rational point gives an eventually defined sequence of constant distance
from the basepoint. In a pointed ultralimit these points form an isometric
copy of $\mathbb Q^a$. Completion extends it to $\mathbb R^a$, preserving
the joint prices on each finite placement.
\end{proof}

By Lemma~\ref{lem:transfer}, failure of Theorem~\ref{thm:flats} gives a
finite prime field $F$ over which some dimension is missing. Realized
dimensions are downward closed. To obtain a spectral line, start with
the two-party flattening rank and insert a unit segment by
Corollary~\ref{cor:spectral-geodesic}. Rescale the segment to lengths
tending to infinity and take a pointed ultralimit based at its midpoint.
There is consequently a largest realized dimension $k\ge1$. We fix
$F$ and $k$ through the contradiction in Section~\ref{sec:diagonal}.
Constants in that argument may depend on both.

Consider the distance matrix of the corners of the Euclidean unit
$(k+1)$-cube. In maximum norm this matrix has positive distance from
the set of distance matrices of that size realizable in complete
$\operatorname{CAT}(0)$ spectral spaces over $F$, allowing all positive
rescalings. Otherwise
a pointed ultralimit would realize the corners exactly. The flat
rectangle theorem then fills in the cube by successive interpolation
on its coordinate edges (Section~\ref{sec:crossing}). Rescaling about
its center by factors tending to infinity and taking a pointed ultralimit
would produce an entire $(k+1)$-flat.
